\section*{Chapitre \ref{ch:qm:girsanov-and-changes-of-numeraire} --- Girsanov et changements de numéraire}

\begin{solution}{exo:qm:girsanov-and-changes-of-numeraire:1}
$W_T = W^{\mathbb Q}_T + \gamma T$ avec $W^{\mathbb Q}$ un $\mathbb Q$-mouvement brownien~: $W_T \sim \mathcal N(2, 4)$ sous $\mathbb Q$.
\end{solution}

\begin{solution}{exo:qm:girsanov-and-changes-of-numeraire:2}
Dérive de 4~\%~; $\gamma = -(\mu - r)/\sigma = -0.05/0.25 = -0.2$, l'opposé du ratio de Sharpe.
\end{solution}

\begin{solution}{exo:qm:girsanov-and-changes-of-numeraire:3}
$P(0,5)\,\E^5[X] = 0.8343 \times 2.4 = 2.002$.
\end{solution}

\begin{solution}{exo:qm:girsanov-and-changes-of-numeraire:4}
Sous $\mathbb Q^T$, la dérive de $r$ est diminuée de $\sigma^2B(T - t) > 0$~: les états à taux élevés, dans lesquels l'obligation est bon marché,
reçoivent moins de poids. $B(5) = 1.8358$ et $\sigma^2B^2/2 = 0.0001 \times 3.370/2 = 1.7$ pb, l'écart entre
3.918~\% et 3.901~\%.
\end{solution}

\begin{solution}{exo:qm:girsanov-and-changes-of-numeraire:5}
$S_t = (P(t,T_0) - P(t,T_n))/A_t$ est un portefeuille négocié divisé par le \omterm{def:qm:girsanov-and-changes-of-numeraire:numeraire}{numéraire} $A_t$, donc une
$\mathbb Q^A$-martingale. Une swaption receveuse verse $A_{T_0}(K - S_{T_0})^+$ en $T_0$, donc sa valeur est égale à $A_0\,\E^A[(K -
S_{T_0})^+]$.
\end{solution}

\begin{solution}{exo:qm:girsanov-and-changes-of-numeraire:6}
Sous $\mathbb Q^{S^2}$, le rapport $S^1/S^2$ est une \omterm{def:qm:probability-at-speed:martingale}{martingale} de volatilité $\sqrt{0.2^2 + 0.2^2 - 2 \times 0.5
\times 0.04} = 0.2$, et la valeur est $S^2_0\,\E^{S^2}[(S^1_T/S^2_T - 1)^+] = 100(\Phi(0.1) - \Phi(-0.1)) = 7.97$~:
formule de Margrabe, sans aucun taux d'intérêt.
\end{solution}

\begin{solution}{exo:qm:girsanov-and-changes-of-numeraire:7}
Soit $X = \mu t + \sigma W$. Sous $d\mathbb Q/d\P = e^{-(\mu/\sigma)W_T - \mu^2T/2\sigma^2}$, $X/\sigma$ est un \omterm{def:qm:brownian-motion:bm}{mouvement
brownien} standard, et $d\P/d\mathbb Q = e^{(\mu/\sigma^2)X_T - \mu^2T/2\sigma^2}$ ne dépend de la trajectoire qu'à travers $X_T$. Avec $m =
\min_{t\le T}X_t$ et $b < 0$, le principe de réflexion donne, sous $\mathbb Q$, la loi conjointe $\mathbb Q(m \le b, X_T \in dx)
= \mathbb Q(X_T \in 2b - dx)$ pour $x \ge b$. D'où
\[
\P(m \le b) = \P(X_T \le b) + \int_b^\infty e^{\mu x/\sigma^2 - \mu^2T/2\sigma^2}\,\varphi_{\sigma\sqrt T}(2b - x)\,dx,
\]
où $\varphi_s$ est la densité de la loi $\mathcal N(0, s^2)$~; en complétant le carré, l'intégrale vaut $e^{2\mu b/\sigma^2}\Phi((b +
\mu T)/\sigma\sqrt T)$, et $\P(X_T \le b) = \Phi((b - \mu T)/\sigma\sqrt T)$. Pour l'arrêt du chapitre~2 ($b = \ln
0.98$, $\mu = -0.045$, $\sigma = 0.30$, $T = 21/252$)~: 82.4~\%.
\end{solution}

\begin{solution}{exo:qm:girsanov-and-changes-of-numeraire:8}
Un prix est une espérance sous une probabilité \omterm{def:qm:probability-at-speed:martingale}{martingale} pour le \omterm{def:qm:girsanov-and-changes-of-numeraire:numeraire}{numéraire} choisi, et non sous la
probabilité historique~: avec le \omterm{def:qm:girsanov-and-changes-of-numeraire:numeraire}{compte monétaire} comme \omterm{def:qm:girsanov-and-changes-of-numeraire:numeraire}{numéraire}, l'action doit dériver au taux
court. Utiliser la dérive prévue revient à coter une créance différente, susceptible d'arbitrage~; la prévision relève de la
décision de trading (acheter si le prix de marché est inférieur à la valeur qu'implique sa propre vision), et non dans la
probabilité de valorisation.
\end{solution}

\begin{solution}{pb:qm:girsanov-and-changes-of-numeraire:1}
\textbf{1.} $P(0,5) = 0.8343$, $P(0,5.25) = 0.8262$.
\textbf{2.} $(0.8343/0.8262 - 1)/0.25 = 3.92\%$.
\textbf{3.} 3.918~\% et 3.901~\%~: sous la \omterm{def:qm:girsanov-and-changes-of-numeraire:forward}{probabilité forward}, la dérive de $r$ est plus faible de $\sigma^2B(5 - t)$.
\textbf{4.} Un put de strike $x = 1/(1 + \delta K) = 0.98888$ sur $1 + \delta K = 1.01125$ obligations zéro-coupon venant à échéance
en 5.25, expirant à 5.
\textbf{5.} $\sigma_P = 0.01\sqrt{(1 - e^{-5})/1}\times B(0.25) = 0.234\%$.
\textbf{6.} 3.26~pb de notionnel~: 32\,600~USD.
\textbf{7.} 260 (pas hebdomadaires sur cinq ans) contre un.
\textbf{8.} $3.34 \pm 0.05$ pb.
\textbf{9.} $3.32 \pm 0.05$ pb.
\textbf{10.} Oui~: à 1.6 et 1.2 erreurs types de 3.26.
\textbf{11.} $(P(5,5)/P(0,5))/(B_5/B_0) = 1/(P(0,5)B_5)$, avec $B_5 = e^{\int_0^5r}$ le long de la trajectoire.
\textbf{12.} $3.906\% \pm 0.003\%$, contre $f(0,5) = 3.901\%$~; les pondérations valent en moyenne 1.000.
\textbf{13.} 0.996~: l'aléa du facteur d'actualisation est faible devant celui du payoff, donc le supprimer
ne réduit guère la variance par trajectoire.
\textbf{14.} Le coût par trajectoire~: une normale au lieu de 260, donc environ 260 fois moins de calcul pour la même
erreur, puis la formule fermée, sans aucun tirage.
\textbf{15.} Non~: $3.32 \pm 0.05$ pb avec 208 pas par an, dans le bruit de 3.34.
\textbf{16.} Chaque caplet sous la \omterm{def:qm:girsanov-and-changes-of-numeraire:forward}{probabilité forward} de sa propre date, puis on somme~; une simulation conjointe de toutes
les dates exige une seule probabilité pour toutes, la probabilité spot ou terminale du livre~6.
\textbf{17.} La \omterm{def:qm:girsanov-and-changes-of-numeraire:annuity}{probabilité annuité}~; le taux de swap au pair en est la \omterm{def:qm:probability-at-speed:martingale}{martingale}.
\textbf{18.} Que chaque prix simulé divisé par le \omterm{def:qm:girsanov-and-changes-of-numeraire:numeraire}{numéraire} a une moyenne constante d'une date à l'autre (le
test de \omterm{def:qm:probability-at-speed:martingale}{martingale} déflatée), et que les pondérations du \omterm{def:qm:girsanov-and-changes-of-numeraire:changenum}{changement de numéraire} ont une moyenne égale à un.
\textbf{19.} Résultat nommé~: \emph{deux probabilités, un seul prix}~: le caplet vaut 3.26~pb de notionnel
(32\,600~USD sur 100~millions de dollars)~; les simulations en \omterm{def:qm:girsanov-and-changes-of-numeraire:emm}{probabilité risque-neutre} et en \omterm{def:qm:girsanov-and-changes-of-numeraire:forward}{probabilité forward} s'accordent avec lui et ont
la même erreur type par trajectoire, mais la \omterm{def:qm:girsanov-and-changes-of-numeraire:forward}{probabilité forward} demande une normale par trajectoire au lieu de 260,
et donne la formule fermée.
\textbf{20.} Le \omterm{def:qm:girsanov-and-changes-of-numeraire:changenum}{changement de numéraire} modifie la dérive et l'unité dans laquelle les valeurs sont comptées, jamais les prix ni
la volatilité.
\end{solution}

\begin{solution}{iq:qm:girsanov-and-changes-of-numeraire:1}
Il ajoute une dérive~: $W - \int\gamma\,dt$ est un \omterm{def:qm:brownian-motion:bm}{mouvement brownien} sous la nouvelle probabilité. Il ne peut pas modifier la
\omterm{def:qm:brownian-motion:qv}{variation quadratique}, qui est visible sur une seule trajectoire, donc les volatilités sont les mêmes sous toutes les
\omterm{def:qm:probability-at-speed:change}{mesures équivalentes}.
\iqlookfor{dérive oui, volatilité non, et pourquoi.}
\end{solution}

\begin{solution}{iq:qm:girsanov-and-changes-of-numeraire:2}
La \omterm{def:qm:girsanov-and-changes-of-numeraire:emm}{probabilité risque-neutre} est un outil de valorisation, non une prévision~: c'est la probabilité sous laquelle les prix
exprimés en unités du \omterm{def:qm:girsanov-and-changes-of-numeraire:numeraire}{compte monétaire} forment des jeux équitables. Une action détenue face à une position courte
en \omterm{def:qm:girsanov-and-changes-of-numeraire:numeraire}{compte monétaire} doit alors ne rien rapporter en moyenne dans ces unités, donc sa dérive vaut $r$~; le
rendement espéré réel se situe sous $\P$.
\iqlookfor{probabilité comme unité de compte, non comme croyance.}
\end{solution}

\begin{solution}{iq:qm:girsanov-and-changes-of-numeraire:3}
La \omterm{def:qm:girsanov-and-changes-of-numeraire:emm}{probabilité martingale équivalente} pour l'obligation zéro-coupon échéant en $T$. Le taux forward simple pour
$[S, T]$ est $(P(t,S)/P(t,T) - 1)/\delta$~: un prix négocié divisé par le \omterm{def:qm:girsanov-and-changes-of-numeraire:numeraire}{numéraire}, moins une constante, donc une
\omterm{def:qm:probability-at-speed:martingale}{martingale}.
\iqlookfor{\omterm{def:qm:girsanov-and-changes-of-numeraire:numeraire}{numéraire} $P(t,T)$ et argument du rapport.}
\end{solution}

\begin{solution}{iq:qm:girsanov-and-changes-of-numeraire:4}
Avec l'actif 2 comme \omterm{def:qm:girsanov-and-changes-of-numeraire:numeraire}{numéraire}, le payoff $(S^1_T - S^2_T)^+ = S^2_T(S^1_T/S^2_T - 1)^+$ devient un call sur le
\omterm{def:qm:probability-at-speed:martingale}{martingale} $S^1/S^2$ de strike un, de volatilité $\sqrt{\sigma_1^2 + \sigma_2^2 - 2\rho\sigma_1\sigma_2}$~: la formule de Margrabe
indépendante des taux.
\iqlookfor{choisir le \omterm{def:qm:girsanov-and-changes-of-numeraire:numeraire}{numéraire} qui élimine un actif.}
\end{solution}

\begin{solution}{iq:qm:girsanov-and-changes-of-numeraire:5}
Simuler des instruments négociés dont les prix sont connus (obligations zéro-coupon, forwards) et vérifier que
leurs prix divisés par le \omterm{def:qm:girsanov-and-changes-of-numeraire:numeraire}{numéraire} du moteur ont des moyennes constantes à toutes les dates, à quelques
erreurs types près~; vérifier la parité call-put sur les payoffs simulés~; et valoriser la même créance sous deux
\omterm{def:qm:girsanov-and-changes-of-numeraire:numeraire}{numéraires} avec la repondération.
\iqlookfor{tests de \omterm{def:qm:probability-at-speed:martingale}{martingale} sur des instruments dont les prix sont connus.}
\end{solution}

\begin{solution}{iq:qm:girsanov-and-changes-of-numeraire:6}
L'échantillonnage direct ne voit un dépassement qu'une fois sur 3.5 millions de tirages. Tirer $Y \sim \mathcal N(5, 1)$ à la place et moyenner
$e^{-5Y + 12.5}\mathbf 1_{\{Y > 5\}}$~: 10\,000 tirages donnent $2.79 \times 10^{-7}$ avec une erreur type de
$0.07 \times 10^{-7}$, contre la valeur exacte $2.87 \times 10^{-7}$.
\iqlookfor{échantillonnage préférentiel par translation de la moyenne, avec le poids du rapport de vraisemblance.}
\end{solution}
